\documentclass[11pt]{article}
\usepackage[a4paper,margin=18mm]{geometry}
\usepackage{fontspec}
\setmainfont{Latin Modern Roman}
\usepackage{amsmath,amssymb,amsthm,amscd,graphicx,color}
\usepackage{xurl}
\usepackage[unicode,hidelinks]{hyperref}
\allowdisplaybreaks
\setlength\emergencystretch{3em}

\makeatletter
\newtheorem{theorem}{Theorem}
\newtheorem*{theorem*}{Theorem}
\newtheorem{lemma}{Lemma}
\newtheorem*{lemma*}{Lemma}
\newtheorem{corollary}{Corollary}
\newtheorem*{corollary*}{Corollary}
\newtheorem{proposition}{Proposition}
\newtheorem*{proposition*}{Proposition}
\newtheorem{conjecture}{Conjecture}
\newtheorem*{conjecture*}{Conjecture}
{\theoremstyle{definition} \newtheorem{definition}{Definition}
\newtheorem*{definition*}{Definition}
\newtheorem{example}{Example}
\newtheorem*{example*}{Example}
\newtheorem{remark}{Remark}
\newtheorem*{remark*}{Remark}
\newtheorem{note}{Note}
\newtheorem*{note*}{Note}
}

\makeatother
\numberwithin{equation}{section}

\numberwithin{theorem}{section}
\numberwithin{proposition}{section}
\numberwithin{lemma}{section}
\numberwithin{corollary}{section}
\numberwithin{definition}{section}
\numberwithin{example}{section}
\numberwithin{remark}{section}
\numberwithin{note}{section}

\usepackage[all]{xy}

\theoremstyle{remark}
\newtheorem{rem}{Remark}[section]\newtheorem*{notation}{Notation}

\renewcommand{\theequation}{\thesection.\arabic{equation}}

\newcommand{\thmref}[1]{Theorem~\ref{#1}}
\newcommand{\secref}[1]{\S\ref{#1}}
\newcommand{\lemref}[1]{Lemma~\ref{#1}}

\newcommand{\mR}{\mathbb{R}}
\newcommand{\mC}{\mathbb{C}}
\newcommand{\mN}{\mathbb{N}}
\newcommand{\mE}{\mathbb{E}}
\newcommand{\mZ}{\mathbb{Z}}
\newcommand{\mS}{\mathbb{S}}
\newcommand{\mV}{\mathbb{V}}

\newcommand{\cD}{\mathcal{D}}
\newcommand{\cM}{\mathcal{M}}
\newcommand{\cH}{\mathcal{H}}
\newcommand{\cE}{\mathcal{E}}
\newcommand{\cF}{\mathcal{F}}
\newcommand{\cP}{\mathcal{P}}
\newcommand{\cC}{\mathcal{C}}
\newcommand{\cR}{\mathcal{R}}
\newcommand{\cK}{\mathcal{K}}
\newcommand{\cS}{\mathcal{S}}
\newcommand{\cG}{\mathcal{G}}
\newcommand{\cL}{\mathcal{L}}

\newcommand{\dD}{\textbf{D}}

\newcommand{\ux}{\underline{x}}
\newcommand{\uz}{\underline{z}}

\newcommand{\uxb}{\underline{x} \grave{}}
\newcommand{\uxi}{\underline{\xi}}
\newcommand{\uyb}{\underline{y} \grave{}}
\newcommand{\uy}{\underline{y}}
\newcommand{\uub}{\underline{u} \grave{}}
\newcommand{\uu}{\underline{u}}
\newcommand{\uvb}{\underline{v} \grave{}}
\newcommand{\uv}{\underline{v}}

\newcommand{\pj}{\partial_{x_j}}
\newcommand{\pjb}{\partial_{{x \grave{}}_{j}}}
\newcommand{\pkb}{\partial_{{x \grave{}}_{k}}}

\newcommand{\pI}{\partial_{x_i}}
\newcommand{\pIb}{\partial_{{x \grave{}}_{i}}}
\newcommand{\pIcont}{\partial_{x_i} \rfloor}
\newcommand{\pIbcont}{\partial_{{x \grave{}}_{i}} \rfloor}

\newcommand{\pk}{\partial_{x_k}}
\newcommand{\px}{\partial_x}
\newcommand{\py}{\partial_y}
\newcommand{\upx}{\partial_{\underline{x}}}
\newcommand{\upxb}{\partial_{\underline{{x \grave{}}} }}
\newcommand{\upy}{\partial_{\underline{y}}}
\newcommand{\upyb}{\partial_{\underline{{y \grave{}}} }}

\newcommand{\upxa}{\partial_{\underline{x},a}}
\newcommand{\gog}{{\mathfrak g}}
\newcommand{\gom}{{\mathfrak m}}
\newcommand{\gon}{{\mathfrak n}}
\newcommand{\gop}{{\mathfrak p}}

\def\a{{\alpha}}
\def\b{{\beta}}
\def\g{{\gamma}}
\def\k{{\kappa}}
\def\t{{\theta}}
\def\l{{\lambda}}
\def\d{{\delta}}
\def\o{{\omega}}
\def\s{{\sigma}}
\def\la{{\langle}}
\def\ra{{\rangle}}


\makeatletter\newlabel{section8}{{8}{}{Original source locator}{}{}}\makeatother
\begin{document}
\begin{center}\Large Convergent Clifford-valued Hermite-semigroup kernel series and equality on the explicit Laguerre-monogenic basis for kappa=\allowbreak{}0, c>-1, dimension m>=\allowbreak{}2 and Re omega>0\end{center}
\noindent\textbf{Stable ID:} 2029\par
\noindent\textbf{Authors:} Hendrik De Bie; Bent Orsted; Petr Somberg; Vladimir Soucek\par
\noindent\textbf{Original work:} The Clifford Deformation of the Hermite Semigroup\par
\noindent\textbf{Version:} Official SIGMA 9 (2013) article 010, DOI 10.3842/SIGMA.2013.010; journal PDF and native TeX acquired 2026-10-01, exact bytes pinned by SHA256\par
\noindent\textbf{Selected task scope:} Full Theorem6.2 for the undeformed Dunkl case kappa=\allowbreak{}0,dimensionm>=\allowbreak{}2,c>-1,Re omega>0. The original Clifford-valued A,B Bessel-Gegenbauer series converge and its weighted integral transform agrees with the specified spectral semigroup on every basis vector phi\_(t,l,index). No claim of pointwise integration for every L2 representative or boundary Re omega=\allowbreak{}0 is added.\par
\noindent\textbf{Difficulty:} H2: The proof constructs separate reproducing kernels for monogenic and vector-monogenic sectors and matches their radial Bessel-Laguerre integrals with the deformed oscillator eigenvalues. It then controls all angular Gegenbauer and growing-order Bessel contributions by a gamma-denominator estimate to obtain the convergent kernel series, beyond formal spectral exponentiation or a classical Mehler substitution.\par
\noindent\textbf{Editorial boundary note (not author text):} Full original Theorem 6.2, explicit convergent A/B Clifford-valued kernel series and equality on the exact Laguerre-monogenic basis are retained with all new Funk-Hecke/reproducing-kernel, spectral-matching and Gamma-dominant series estimates. The selected branch is kappa=\allowbreak{}0, real c>-1, dimension m>=\allowbreak{}2 and Re omega>0, with the stated weighted Hilbert space and h(r) measure. The m=\allowbreak{}2 limiting-Gegenbauer convention is supplied. Earlier osp(1|2)/basis action, standard Bessel/Laguerre integral and Appendix A bounds remain named prior background with their exact definitions. Original intermediate measure notation is unchanged with its notice; the final weighted transform is the selected statement. No unitary boundary limit, Fourier variant, uncertainty corollary or analytic continuation through Re omega=\allowbreak{}0 is selected.\par
\noindent\textbf{Source:} \url{https://sigma-journal.com/2013/010/}\quad DOI: \url{https://doi.org/10.3842/SIGMA.2013.010}\par
\noindent\textbf{License and attribution:} Copyright retained by the named authors. Published by SIGMA under CC BY 4.0: \url{https://creativecommons.org/licenses/by/4.0/}. Source: \url{https://sigma-journal.com/about.html}. Adaptation consists of standalone excerpt formatting, cover and numbering restoration; original author language and mathematical text are retained.\par
\noindent\textbf{Editorial symbol notice (not author text):} Several intermediate matching integrals in Section6.1 write dx without the weight, whereas the initial ansatz and the selected Theorem6.2 explicitly use h(r)dx and the displayed radial exponent is provided. Exact source pages retained; the final weighted transform is the selected statement.\par
\noindent\textbf{External original locator section8:} Section 8 in the original article discusses additional kernel properties. This is an incidental limitation pointer; the selected non-Dunkl semigroup kernel construction and convergence proof are fully retained.\par
\medskip\hrule\medskip\noindent\textbf{Author text begins below.}\par
\setcounter{section}{1}
\setcounter{subsection}{0}
\setcounter{subsubsection}{0}
\setcounter{equation}{0}
% Original source lines 250--1069
\section{Preliminaries}\label{section2}

In this section we collect some basic results on Clif\/ford algebras and Dunkl operators.

\subsection{Clif\/ford algebras}\label{section2.1}
Let $\mV$ be a~vector space of dimension $m$ with a~given negative def\/inite quadratic form and
let~$\cC l_{m}$ be the corresponding Clif\/ford algebra.
If $\{e_i\}$ is an orthonormal basis of $\mV$, then $\cC l_{m}$ is generated by $e_{i}$, $i=1,
\ldots,m$, with the relations
\begin{gather*} %\label{eq:eij}
e_{i} e_{j} + e_{i} e_{j} = 0, \quad i \neq j,\qquad
 e_{i}^{2} = -1.
\end{gather*}
The algebra $\cC l_{m}$ has dimension $2^{m}$ as a~vector space over $\mR$.
It can be decomposed as $\cC l_{m}=\oplus_{k=0}^{m}\cC l_{m}^{k}$
with $\cC l_{m}^{k}$ the space of $k$-vectors def\/ined by
\begin{gather*}
\cC l_{m}^{k}:=\text{span}\{e_{i_{1}}\cdots e_{i_{k}},\; i_{1}<\cdots<i_{k}\}.
\end{gather*}
The projection on the space of $k$-vectors is denoted by $[\cdot]_{k}$.

The operator $\bar{.}$ is the main anti-involution on the Clif\/ford algebra $\cC l_{m}$ def\/ined by
\begin{gather*}
\overline{a b}=\overline{b}\overline{a},\qquad\overline{e_{i}}=-e_{i},\quad i=1,\ldots,m.
\end{gather*}
Similarly we have the automorphism $\epsilon$ given by
\begin{gather*}
\epsilon(a b)=\epsilon(a)\epsilon(b),\qquad\epsilon(e_{i})=-e_{i},\quad i=1,\ldots,m.
\end{gather*}

In the sequel, we will always consider functions $f$ taking values in $\cC l_{m}$, unless
explicitly mentioned.
Such functions can be decomposed as
\begin{gather*}%\label{clifford_func}
f(x)=f_{0}(x)+\sum_{i=1}^{m}e_{i}f_{i}(x)+\sum_{i<j}e_{i}e_{j}f_{ij}(x)+\cdots+e_{1}
\cdots e_{m}f_{1\ldots m}(x)
\end{gather*}
with $f_{0},f_{i},f_{ij},\ldots,f_{1\ldots m}$ all real-valued functions.

Several important groups can be embedded in the Clif\/ford algebra.
Note that the space of $1$-vectors in $\cC l_{m}$ is canonically isomorphic to $\mV\cong\mR^{m}$.
Hence we can def\/ine
\begin{gather*}
\mathrm{Pin}(m)=\left\{s_{1}s_{2}\cdots s_{n}\,|\,n\in\mN,s_{i}\in\cC l_{m}^{1}\ \text{such that}\ s_{i}^{2}=-1\right\},
\end{gather*}
i.e., the Pin group is the group of products of unit vectors in~$\cC l_{m}$.
This group is a~double cover of the orthogonal group~${\rm O}(m)$ with covering map $p:\mathrm{Pin}(m)
\rightarrow {\rm O}(m)$, which we will describe explicitly in the next section.

Similarly we def\/ine
\begin{gather*}
\mathrm{Spin}(m)=\left\{s_{1}s_{2}\cdots s_{2n}\,|\,n\in\mN,s_{i}\in\cC l_{m}^{1} \ \text{such that} \ s_{i}^{2}=-1\right\},
\end{gather*}
i.e., the Spin group is the group of even products of unit vectors in $\cC l_{m}$.
This group is a~double cover of ${\rm SO}(m)$.
For more information about Clif\/ford algebras and analysis, we refer the reader to~\cite{MR1169463,MR1130821}.

\subsection{Dunkl operators}\label{section2.2}
Denote by $\langle\cdot,\cdot
\rangle$ the standard Euclidean scalar product in $\mR^{m}$ and by $|x|=\langle x,
x\rangle^{1/2}$ the associated norm.
For $\alpha\in\mR^{m}\setminus \{0\}$, the ref\/lection $r_{\alpha}$ in the hyperplane orthogonal to
$\alpha$ is given by
\begin{gather*}
r_{\alpha}(x)=x-2\frac{\langle\alpha,x\rangle}{|\alpha|^{2}}\alpha,\qquad x\in\mR^{m}.
\end{gather*}

A root system is a~f\/inite subset $R\subset\mR^{m}$ of non-zero vectors such that, for every
$\alpha\in R$, the associated ref\/lection $r_{\alpha}$ preserves $R$.
We will assume that $R$ is reduced, i.e.
$R\cap\mR\alpha=\{\pm\alpha\}$ for all $\alpha\in R$.
Each root system can be written as a~disjoint union $R=R_{+}\cup(-R_{+})$, where $R_{+}$ and~$-R_{+}$ are separated by a~hyperplane through the origin.
$R_+$ is called a~positive subsystem of the root system $R$.
The subgroup $\cG\subset {\rm O}(m)$ generated by the ref\/lections $\{r_{\alpha}|\alpha\in R\}$ is
called the f\/inite ref\/lection group associated with~$R$.
We will also assume that $R$ is normalized such that $\langle\alpha,\alpha\rangle=2$ for all
$\alpha\in R$.
For more information on f\/inite ref\/lection groups we refer the reader to~\cite{Humph}.

If we identify $\alpha$ with a~$1$-vector in $\cC l_{m}$ (and hence $\alpha/\sqrt{2}$ with an
element in $\mathrm{Pin}(m)$), we can rewrite the ref\/lection $r_{\alpha}$ as
\begin{gather*}
r_{\alpha}(x)=\frac{1}{2}\alpha\ux\alpha
\end{gather*}
with $\ux=\sum\limits_{i=1}^{m}e_{i}x_{i}$.
Generalizing this map gives us the covering map $p$ from $\mathrm{Pin}(m)$ to ${\rm O}(m)$ as
\begin{gather*}
p(s)(x)=\epsilon(s)\ux s^{-1},\qquad s\in \mathrm{Pin}(m).
\end{gather*}
In particular, we obtain a~double cover of the ref\/lection group $\cG$ as $\widetilde{\cG}=
p^{-1}(\cG)$ (see also the discussion in~\cite{BCT}).

A multiplicity function $\kappa$ on the root system $R$ is a~$\cG$-invariant function $\kappa:R
\rightarrow\mC$, i.e.
$\k(\alpha)=\k(h\alpha)$ for all $h\in\cG$.
We will denote $\k(\alpha)$ by $\k_{\alpha}$.
We will always assume that the multiplicity function is real and satisf\/ies $\k\geq0$.
This assumption is, e.g., 
necessary to obtain the subsequent formula~\eqref{skewnessDunkl}, which is crucial for the sequel.

Fixing a~positive subsystem $R_{+}$ of the root system~$R$ and a~multiplicity function~$\k$, we
introduce the Dunkl operators~$T_{i}$ associated to~$R_{+}$ and $\k$ by (see~\cite{MR951883,
MR1827871})
\begin{gather*}
T_{i}f(x)=\partial_{x_{i}}f(x)+\sum_{\alpha\in R_{+}}\k_{\alpha}\alpha_{i}\frac{f(x)-
f(r_{\alpha}(x))}{\langle\alpha,x\rangle},\qquad f\in C^{1}(\mR^{m}).
\end{gather*}
An important property of the Dunkl operators is that they commute, i.e.
$T_{i}T_{j}=T_{j}T_{i}$.

The Dunkl Laplacian is given by $\Delta_{\k}=\sum\limits_{i=1}^{m}T_i^2$, or more explicitly by
\begin{gather*}
\Delta_{\k}f(x)=\Delta f(x)+2\sum_{\alpha\in R_{+}}\k_{\alpha}\left(\frac{\langle\nabla
f(x),\alpha\rangle}{\langle\alpha,x\rangle}-\frac{f(x)-f(r_{\alpha}(x))}{\langle\alpha,
x\rangle^{2}}\right)
\end{gather*}
with $\Delta$ the classical Laplacian and $\nabla$ the gradient operator.
We also def\/ine the constant
\begin{gather*}
\mu=\frac{1}{2}\Delta_{\k}|x|^2=m+2\sum_{\alpha\in R_+}\k_{\alpha},
\end{gather*}
called the Dunkl-dimension.

It is possible to construct an intertwining operator $V_{\k}$ connecting the classical derivatives
$\partial_{x_{j}}$ with the Dunkl operators $T_{j}$ such that $T_{j}V_{\k}=V_{\k}
\partial_{x_{j}}$ (see, e.g.,~\cite{MR1273532}).
Note that explicit formulae for~$V_{\k}$ are only known in a~few special cases.

The weight function related to the root system $R$ and the multiplicity function $\k$ is given by
$w_{\k}(x)=\prod\limits_{\alpha\in R_{+}}|\langle\alpha,x\rangle|^{2\k_{\alpha}}$.
For suitably chosen functions $f$ and $g$ one then has the following property of integration by
parts (see~\cite{Du4})
\begin{gather}
\int_{\mR^{m}}(T_{i}f)g w_{\k}(x)dx=-\int_{\mR^{m}}f\left(T_{i}g\right)w_{\k}(x)dx.
\label{skewnessDunkl}
\end{gather}
For more information about Dunkl operators we refer the reader to~\cite{MR1827871, MR2022853}.

The starting point in the subsequent analysis is the Dunkl--Dirac operator, given by
\begin{gather*}
\cD_{\k}=\sum_{i=1}^{m}e_{i}T_{i}.
\end{gather*}
Together with the vector variable $\ux=\sum\limits_{i=1}^{m}e_{i}x_{i}$ this Dunkl--Dirac operator
generates a~copy of~$\mathfrak{osp}(1|2)$, see~\cite{Orsted} or the subsequent Theorem~\ref{ospFamily}.
In particular, we have
\begin{gather*}
\cD_{\k}^{2}=-\Delta_{\k}\qquad\text{and}\qquad\ux^{2}=-|\ux|^{2}=-r^{2}=-
\sum_{i=1}^{m}x_{i}^{2}.
\end{gather*}

\section{Intertwining operators}\label{section3}

Let, for $a,b\in\mR$, $P$ and $Q$ be two operators def\/ined by
\begin{gather*}
P f(\ux)=r^{b}f\left(\left(\frac{a}{2}\right)^{\frac{1}{a}}\ux r^{\frac{2}{a}-1}\right),\qquad
Q f(\ux)=r^{-\frac{ab}{2}}f\left(\left(\frac{2}{a}\right)^{\frac{1}{2}}\ux r^{\frac{a}{2}-1}\right).
\end{gather*}
These two operators act as generalized Kelvin transformations.
Indeed, one can easily compute their composition
\begin{gather*}
Q P=P Q=\left(\frac{2}{a}\right)^{\frac{b}{2}}.
\end{gather*}
We will show that these operators allow to reduce the Dirac operator $\dD$ to a~simpler form.

We have the following proposition, where $\mE=\sum\limits_{i=1}^{m}x_{i}\partial_{x_{i}}$ denotes the
Euler operator.
Recall also from the introduction that $\ux_{a}=r^{\frac{a}{2}-1}\ux$.

\begin{proposition}
One has the following intertwining relations
\begin{gather*}
\left(\frac{a}{2} \right)^{\frac{b-1}{2}} Q \left( \cD_{\k} + b r^{-2} \ux + c r^{-2}\ux \mE
\right) P = r^{1- \frac{a}{2}}\cD_{k} + \b r^{-\frac{a}{2}-1} \ux + \g r^{-\frac{a}{2}-1}\ux \mE,\\
\left(\frac{a}{2} \right)^{\frac{b+1}{2}} Q \ux P = \ux_{a}
\end{gather*}
with
$
\b= 2b+ bc$, $
\g= \frac{2}{a} (1+c) -1$.
\end{proposition}

\begin{proof}
In~\cite[Proposition 3]{H12}, we already proved that
\begin{gather*}
\left(\frac{a}{2}\right)^{\frac{b-1}{2}}Q\left(\cD_{\k}\right)P=r^{1-\frac{a}{2}}\cD_{k}+
b r^{-\frac{a}{2}-1}\ux+\left(\frac{2}{a}-1\right)r^{-\frac{a}{2}-1}\ux\mE=\ux_{a}.
\end{gather*}
Similarly we obtain
\begin{gather*}
\left(\frac{a}{2}\right)^{\frac{b-1}{2}}Q\left(r^{-2}\ux\right)P=r^{-\frac{a}{2}-1}\ux
\end{gather*}
and
\begin{gather*}
\left(\frac{a}{2}\right)^{\frac{b-1}{2}}Q\left(r^{-2}\ux\mE\right)P=b r^{-\frac{a}{2}-1}
\ux+\left(\frac{2}{a}\right)r^{-\frac{a}{2}-1}\ux\mE.
\end{gather*}
This completes the proof of the proposition.
\end{proof}

So we are reduced to the study of the operator
\begin{gather*}
\dD=\cD_{\k}+b r^{-2}\ux+c r^{-2}\ux\mE,
\end{gather*}
where $b,c\in\mR$, $c\neq-1$.
Here, the term $b r^{-2}\ux$ can also be removed.
Indeed, we have
\begin{gather*}
r^{-\alpha}\left(\cD_{\k}+b r^{-2}\ux+c r^{-2}\ux\mE\right)r^{\alpha}=\cD_{\k}+c
r^{-2}\ux\mE,
\end{gather*}
when $\alpha=-b/(1+c)$.

As a~result of the previous discussion, we see that it is suf\/f\/icient to study the function theory
for the operator
\begin{gather*}
\dD=\cD_{\k}+c r^{-2}\ux\mE,
\end{gather*}
where we have put $a=2$, $b=0$.
Furthermore, we will restrict ourselves to the case $c>-1$ for reasons that will become clear in
Proposition~\ref{partIntProp}.
Similarly, we no longer need to consider $\ux_{a}$ but can restrict ourselves to $\ux$.
Now we repeat the basic facts concerning this operator we need in the sequel.
All the results are taken from~\cite{H12}, putting $a=2$, $b=0$.

\begin{theorem}\label{ospFamily}
The operators $\dD$ and $\ux$ generate a~Lie superalgebra, isomorphic to $\mathfrak{osp}(1|2)$,
with the following relations
\begin{alignat*}{3}
& \{\ux,\dD\}=-2(1+c)\left(\mE+\frac{\delta}{2}\right),\qquad&&
\left[\mE+\frac{\delta}{2},\dD\right]=-\dD,& \\
& \left[\ux^{2},\dD\right]=2(1+c)\ux,\qquad&&
\left[\mE+\frac{\delta}{2},\ux\right]=\ux,& \\
& \left[\dD^{2},\ux\right]=-2(1+c)\dD,\qquad&&
\left[\mE+\frac{\delta}{2},\dD^{2}\right]=-2\dD^{2},&\\
& \left[\dD^{2},\ux^{2}\right]=4(1+c)^{2}\left(\mE+\frac{\delta}{2}\right),\qquad&&
\left[\mE+\frac{\delta}{2},\ux^{2}\right]=2\ux^{2},&
\end{alignat*}
where $\delta=1+\frac{\mu-1}{1+c}$.
\end{theorem}

Note that the square of $\dD$ is a~complicated operator, given by
\begin{gather*}
\dD^{2}=-\Delta_{\k}-\left(c\mu\right)r^{-1}\partial_{r}
-\left(c^{2}+2c\right)\partial_{r}^{2}+c r^{-2}\sum_{i}x_{i}T_{i}
-c r^{-2}\sum_{i<j}e_{i}e_{j}(x_{i}T_{j}-x_{j}T_{i}).
\end{gather*}
If $\k=0$, the formula for $\dD^{2}$ simplif\/ies a~bit as now $\sum_{i}x_{i}T_{i}=r
\partial_{r}=\mE$.

\begin{remark}
The operator $\dD=\cD_{\k}+c r^{-2}\ux\mE$ is also considered from a~very dif\/ferent
perspective in~\cite{MR2269930} (in the case $\k=0$), where the eigenfunctions of this operator
are studied.
\end{remark}

Let us now discuss the symmetry of the generators of $\mathfrak{osp}(1|2)$.
First we def\/ine the action of the Pin group on $C^{\infty}(\mR^{m})\otimes\cC l_{m}$ for $s\in
\mathrm{Pin}(m)$ as
\begin{gather*}
\rho(s): \ C^{\infty}(\mR^{m}) \otimes \cC l_{m} \rightarrow C^{\infty}(\mR^{m}) \otimes \cC l_{m},\qquad
 f\otimes b \rightarrow f (p\big(s^{-1}\big)x) \otimes s b.
\end{gather*}
We then have
\begin{proposition}\label{equivarianceG}
Let $s\in\widetilde{\cG}$ and define $\mathrm{sgn}(s):=\mathrm{sgn}(p(s))$.
Then one has
\begin{gather*}
\rho(s) \ux = \mathrm{sgn}(s) \ux \rho(s),\qquad
\rho(s) \dD = \mathrm{sgn}(s) \dD \rho(s).
\end{gather*}
\end{proposition}

\begin{proof}
This follows immediately from the def\/inition of $\rho$ and the $\cG$-equivariance of the Dunkl
operators.
\end{proof}

So up to sign, the Dirac operator $\dD$ is $\widetilde{\cG}$-equivariant.
At this point it is interesting to remark that an algebraic analog of the Dunkl--Dirac operator
$\dD$ for
graded af\/f\/ine Hecke algebras is introduced in~\cite{BCT} with the motivation to prove
a version of Vogan's Conjecture for Dirac cohomology.
The formulation is based on a~uniform geometric parametrization of spin
representations of Weyl groups.
This Dirac operator is an algebraic variant of our family deformation of the dif\/ferential Dirac
operator
for special values of the deformation parameters.
Moreover, it satisf\/ies the same symmetry as in Proposition~\ref{equivarianceG}, see~\cite[Lemma~3.4]{BCT}.

There is a~measure naturally associated with~$\dD$.
Indeed, one has

\begin{proposition}\label{partIntProp}
If $c>-1$, then for suitable differentiable functions $f$ and $g$ one has
\begin{gather*}
\int_{\mR^{m}}\overline{(\dD f)}g h(r)w_{\kappa}(x)dx=\int_{\mR^{m}}\overline{f}
(\dD g)h(r)w_{\k}(x)dx
\end{gather*}
with $h(r)=r^{1-\frac{1+\mu c}{1+c}}$, provided the integrals exist.
\end{proposition}
In this proposition, $\bar{\cdot}$ is the main anti-involution on the Clif\/ford algebra $\cC l_{m}$.

\section[Representation space for the deformation family of the Dunkl-Dirac
operator]{Representation space for the deformation family\\ of the Dunkl--Dirac
operator}\label{section4}

The function space we will work with is $\cL^{2}_{\k,c}(\mR^{m})=L^{2}(\mR^{m},h(r)w_{\k}(x)dx)
\otimes\cC l_{m}$.
This space has the following decomposition
\begin{gather*}
\cL^{2}_{\k,c}(\mR^{m})=L^{2}\big(\mR^{+},r^{\frac{\mu-1}{1+c}}dr\big)\otimes
L^{2}(\mS^{m-1},w_{\k}(\xi)d\sigma(\xi))\otimes\cC l_{m},
\end{gather*}
where on the right-hand side the topological completion of the tensor product is understood and
with $d\sigma(\xi)$ the Lebesgue measure on the sphere $\mS^{m-1}$.
The space $L^{2}(\mS^{m-1},w_{\k}(\xi)d\sigma(\xi))\otimes\cC l_{m}$ can be further decomposed
into Dunkl harmonics and subsequently into Dunkl monogenics.
This leads to
\begin{gather*}
L^{2}\big(\mS^{m-1},w_{\k}(\xi)d\sigma(\xi)\big)\otimes\cC l_{m}=\bigoplus_{\ell=0}^{\infty}
\left(\cM_{\ell}\oplus\ux\cM_{\ell}\right)\big|_{\mS^{m-1}},
\end{gather*}
where $\cM_{\ell}=\ker{\cD_{\k}}\cap\left(\cP_{\ell}\otimes\cC l_{m}\right)$ is the space
of Dunkl monogenics of degree $\ell$, with $\cP_{\ell}$ the space of homogeneous polynomials of
degree $\ell$ (see also~\cite{DBGeg} for more details on Dunkl monogenics).

Using this decomposition, we have obtained in~\cite{H12} a~basis for $\cL^{2}_{\k,c}(\mR^{m})$.
This basis is given by the set $\{\phi_{t,\ell,m}\}$ ($t,\ell\in\mN$ and $m=1,\ldots,\dim
\cM_{\ell}$), def\/ined as
\begin{gather*}
\phi_{2t,\ell,m} = 2^{2t} (1+c)^{2t} t! L_{t}^{\frac{\gamma_{\ell}}{2} -1}(r^{2})r^{\beta_{\ell}}
M_{\ell}^{(m)} e^{-r^{2}/2},\\
\phi_{2t+1,\ell,m} = - 2^{2t+1} (1+c)^{2t+1} t! L_{t}^{\frac{\gamma_{\ell}}{2}}(r^{2}) \ux
r^{\beta_{\ell}} M_{\ell}^{(m)} e^{-r^{2}/2}
\end{gather*}
with $L_{\alpha}^{\beta}$ the Laguerre polynomials and
\begin{gather*}
\beta_{\ell} = - \frac{c}{1+c}\ell,\qquad
\gamma_{\ell} = \frac{2}{1+c}\left( \ell + \frac{\mu-2}{2}\right) + \frac{c+2}{1+c},
\end{gather*}
and where $M_{\ell}^{(m)}$ ($m=1,\ldots,\dim\cM_{\ell}$) forms an orthonormal basis of
$\cM_{\ell}$, i.e.
\begin{gather*}
\left[\int_{\mS^{m-1}}\overline{M_{\ell}^{(m_{1})}}(\xi)M_{\ell}^{(m_{2})}(\xi)w_{\k}(\xi)
d\sigma(\xi)\right]_{0}=\delta_{m_{1}m_{2}}
\end{gather*}
with $[\cdot]_0$ the projection on the scalar part of the Clif\/ford algebra.
The dimension of $\cM_{\ell}$ is given by
\begin{gather*}
\dim_{\mR}{\cM_{\ell}} = \dim_{\mR}{\cC l_{m}} \dim_{\mR}{\cP_{\ell}\big(\mR^{m-1}\big)}
=2^{m} \frac{(\ell +m-2)!}{\ell! (m-2)!}
\end{gather*}
with $\cP_{\ell}\big(\mR^{m-1}\big)$ the space of homogeneous polynomials of degree $\ell$ in $m-1$
variables (see~\cite{MR1169463}).

Using formula (4.10) in~\cite{H12} and the proof of Theorem 3 in \cite{H12}, one obtains the
following formulae for the action of $\dD$ and $\ux$ on the generalized Laguerre functions
\begin{gather}
2 \dD \phi_{t,\ell,m}\! = \phi_{t+1,\ell,m}\! + C(t,\ell) \phi_{t-1,\ell,m},\qquad
-2(1+c) \ux \phi_{t,\ell,m}\! = \phi_{t+1,\ell,m}\! - C(t,\ell) \phi_{t-1,\ell,m}\!\!\!\!\label{ActionOnBasis}
\end{gather}
with
\begin{gather*}
C(2t, \ell)= 4 (1+c)^{2} t,\qquad
C(2t+1, \ell)= 2 (1+c)^{2}(\gamma_{\ell}+ 2t).
\end{gather*}
These formulae determine the action of $\mathfrak{osp}(1|2)$ on $\cL^{2}_{\k,c}(\mR^{m})$.
Recall also that the action of $\widetilde{\cG}$ on $\cL^{2}_{\k,c}(\mR^{m})$ is given by $\rho$
(see Section~\ref{section3}).

Subsequently, we can def\/ine a~creation and annihilation operator in this setting by
\begin{gather}\label{CreaAnn}
A^{+} = \dD - (1+c) \ux,\qquad
A^{-} = \dD + (1+c) \ux
\end{gather}
satisfying
\begin{gather*}%\label{CArel}
A^{+} \phi_{t,\ell,m} = \phi_{t+1,\ell,m},\qquad
A^{-} \phi_{t,\ell,m} = C(t, \ell) \phi_{t-1,\ell,m}.
\end{gather*}

Now we introduce the following inner product
\begin{gather*}
\langle f,g\rangle=\left[\int_{\mR^{m}}\overline{f^{c}}g h(r)w_{\k}(x)dx\right]_{0},
\end{gather*}
where $h(r)$ is the measure associated to $\dD$ (see Proposition~\ref{partIntProp}) and $f^{c}$ is
the complex conjugate of $f$.
It is easy to check that this inner product satisf\/ies
\begin{gather}\label{adjointsIP}
\langle \dD f, g \rangle = \la f, \dD g\ra,\qquad
\la \ux f, g \ra = -\la f, \ux g \ra.
\end{gather}
The related norm is def\/ined by $||f||^{2}=\la f,f\ra$.

\begin{theorem}%\label{orthLaguerre}
We have
\begin{gather*}
\langle\phi_{t_{1},\ell_{1},m_{1}},\phi_{t_{2},\ell_{2},m_{2}}\rangle=c(t_{1},\ell_{1})
\delta_{t_{1}t_{2}}\delta_{\ell_{1}\ell_{2}}\delta_{m_{1}m_{2}},
\end{gather*}
where $c(t,\ell)$ is a~constant depending on $t$ and $\ell$.
\end{theorem}

The functions $\phi_{t,\ell,m}$ are eigenfunctions of the Hamiltonian of a~generalized harmonic
oscillator.
\begin{theorem}
\label{HarmOsc}
The functions $\phi_{t,\ell,m}$ satisfy the following second-order PDE
\begin{gather*}
\left(\dD^{2}-(1+c)^{2}\ux^{2}\right)\phi_{t,\ell,m}=(1+c)^{2}(\gamma_{\ell}+2t)
\phi_{t,\ell,m}.
\end{gather*}
\end{theorem}

\begin{proof}
This follows immediately from the formula~\eqref{ActionOnBasis}.
\end{proof}

Theorem~\ref{HarmOsc} combined with the def\/inition of $A^{+}$, $A^{-}$ in~\eqref{CreaAnn} allows us to
decompose the space $\cL^{2}_{\k,c}(\mR^{m})$ under the action of $\mathfrak{osp}(1|2)$.
Clearly the odd elements $A^{+}$ and $A^{-}$ generate
$\mathfrak{osp}(1|2)$ as they are linear combinations of $\dD$ and $\ux$.
Moreover, they act between two basis vectors
$\{\phi_{t,\ell,m}\}$ of $\cL^{2}_{\k,c}(\mR^{m})$, so it is suf\/f\/icient to consider vectors in an
irreducible representation of $\mathfrak{osp}(1|2)$ inside the functional space.
This is achieved as follows~-- for f\/ixed $\ell$ and $m$ each vector $\phi_{0,\ell,m}$ generates the
irreducible representation
\begin{gather*}
\xymatrix{
\phi_{0,\ell,m}\ar@<1ex>[r]^{A^+}\ar@(dl,dr)_{L}&\phi_{1,\ell,m}\ar@(dl,dr)_{L}
\ar@<1ex>[r]^{A^+}\ar@<1ex>[l]^{A^-}&\phi_{2,\ell,m}\ar@(dl,dr)_{L}\ar@<1ex>[r]^{A^+}
\ar@<1ex>[l]^{A^-}&\phi_{3,\ell,m}\ar@(dl,dr)_{L}\ar@<1ex>[r]^{A^+}\ar@<1ex>[l]^{A^-}
&\phi_{4,\ell,m}\ar@(dl,dr)_{L}\ar@<1ex>[r]^{A^+}\ar@<1ex>[l]^{A^-}&\ar@<1ex>[l]^{A^-}\ldots
}
\end{gather*}
where
\begin{gather*}
L=\frac{1}{2}\{A^{+},A^{-}\}=\dD^{2}-(1+c)^{2}\ux^{2}
\end{gather*}
with the action given in Theorem~\ref{HarmOsc}.
In fact this highest weight representation is labeled by~$\ell$ only and we will denote it
$\pi(\ell)$.
In conclusion, we obtain the decomposition of our functional space $\cL^{2}_{\k,c}(\mR^{m})$
into a~discrete direct sum of highest weight (inf\/inite-dimensional)
Harish-Chandra modules for~$\mathfrak{osp}(1|2)$:
\begin{gather*}
\cL^{2}_{\k,c}(\mR^{m})=\bigoplus_{\ell=0}^{\infty}\pi(\ell)\otimes\cM_{\ell}.
\end{gather*}

These results should be compared with Theorem 3.19 and Section~3.6 in~\cite{Orsted2} (where one
uses~$\mathfrak{sl}_{2}$ instead of $\mathfrak{osp}(1|2)$).
Also notice that the claim should be understood as an assertion on the deformation of the Howe dual
pair for $\mathfrak{osp}(1|2)$ inside the Clif\/ford--Weyl algebra on $\mR^{m}$ acting on a~f\/ixed vector
space $\cL^{2}_{\k,c}(\mR^{m})$.

In particular, we have the following result.
Recall that an operator $T$ is essentially selfadjoint on a~Hilbert space $H$ if $T$ is a
symmetric operator with a~dense domain $D(T)\subset H$ such that for a~complete orthogonal set
$\{f_n\}_n$ in $H$ with $f_n\in D(H)$, there exist $\{\mu_n\}_n$ solving $Tf_n=\mu_nf$ for all
$n\in\mN$.
\begin{proposition}
Let $c>-1$ and $\kappa>0$.
The operator $L$ acting on
$\cL^{2}_{\k,c}(\mR^{m})$ is essentially selfadjoint $($i.e.\
symmetric and its closure is a~selfadjoint operator$)$.
Moreover, $L$ has no continuous spectrum and its
discrete spectrum is given by
\begin{gather*}
\operatorname{Spec}(L)=\{2(1+c)\ell+2(1+c)^{2}t+(1+c)(\mu+c)\,|\,\ell,t\in\mN\}.
\end{gather*}
\end{proposition}

Using Theorem~\ref{HarmOsc} we can now def\/ine the holomorphic semigroup for the deformed Dirac
operator by
\begin{gather*}
\cF_{\dD}^{\omega}=e^{\o\left(\frac{1}{2}+\frac{\mu-1}{2(1+c)}\right)}e^{\frac{-\o}{
2(1+c)^{2}}\left(\dD^{2}-(1+c)^{2}\ux^{2}\right)}.
\end{gather*}
Here, $\omega$ takes values in the right half-plane of $\mC$.
The special boundary value $\o=i\pi/2$ corresponds to the Fourier transform.
In that case, we will use the notation $\cF_{\dD}$.
The functions $\phi_{t,\ell,m}$ are eigenfunctions of $\cF_{\dD}^{\omega}$ satisfying
\begin{gather}
\label{eigvalsFourier}
\cF_{\dD}^{\o}(\phi_{t,\ell,m})=e^{-\o t}e^{-\frac{\o\ell}{(1+c)}}\phi_{t,\ell,m}.
\end{gather}
Note that in the special case $\kappa=0$, $c=0$ the operator $\cF_{\dD}^{\omega}$ reduces to the
classical Hermite semigroup (see, e.g.,~\cite{MR0974332}).

\begin{remark}
One can also consider more general deformations of the Dirac operator, by adding suitable odd
powers of $\Gamma=-\ux\cD_{\k}-\mE$ to $\dD$ as follows
\begin{gather*}
\dD=\cD_{\k}+c r^{-2}\ux\mE+\sum_{j=0}^{\ell}c_{j}r^{-1}\left(\Gamma-\frac{\mu-1}{2}
\right)^{2j+1},\qquad c_{j}\in\mR.
\end{gather*}
This does not alter the $\mathfrak{osp}(1|2)$ relations, as $\Gamma-\frac{\mu-1}{2}$
anti-commutes with $\ux$ and has the correct homogeneity.
In particular, $\Gamma-\frac{\mu-1}{2}$ can be seen as the square root of the Casimir of
$\mathfrak{osp}(1|2)$, see~\cite[Example~2 in Section~2.5]{MR1773773}.
\end{remark}

In the sequel of the paper, we will always assume $\kappa=0$ or in other words, we do not consider
the Dunkl deformation.
This is to simplify the notation of the results.
Most statements can be generalized to the Dunkl case by a~suitable composition with the Dunkl
intertwining operator~$V_{\k}$, except the results obtained in Section~\ref{section8}.

Recall that for $\kappa=0$, the Dunkl--Dirac operator $\cD_{\k}$ reduces to the orthogonal Dirac
operator $\upx=\sum\limits_{i=1}^{m}e_{i}\partial_{x_{i}}$ and the Dunkl dimension~$\mu$ to the ordinary
dimension~$m$.

\section{Reproducing kernels}\label{section5}

In this section we determine the reproducing kernels for $\cM_{k}$ and $\ux\cM_{k}$.
We start with an auxiliary Lemma, which can be thought of as a~Clif\/ford analogue
of the Funk--Hecke transform.
We def\/ine the wedge product of two vectors as
\begin{gather*}
\ux\wedge\uy:=\sum_{j<k}e_{j}e_{k}(x_{j}y_{k}-x_{k}y_{j}).
\end{gather*}

\begin{lemma}
\label{auxintegrals}
Put $\ux=r\ux'$ and $\uy=s\uy'$ with $\ux',\uy'\in\mS^{m-1}$.
Furthermore, put $\lambda=(m-2)/2$ and $\sigma_m=2\pi^{m/2}/\Gamma(m/2)$.
Then one has, with $M_{l}\in\cM_{\ell}$
\begin{gather*}
\int_{\mS^{m-1}} C_{k}^{\lambda}(\la\ux', \uy' \ra) M_{\ell}(\ux')d \s (\ux') = \sigma_m
\frac{\l}{\l+k} \delta_{k, \ell} M_{\ell}(\uy'),\\
\int_{\mS^{m-1}} C_{k}^{\lambda}(\la\ux', \uy' \ra) \ux' M_{\ell}(\ux')d \s (\ux') = \sigma_m
\frac{\l}{\l+k} \delta_{k, \ell+1} \uy' M_{\ell}(\uy'),\\
\int_{\mS^{m-1}} (\ux' \wedge \uy' ) C_{k-1}^{\lambda+1}(\la\ux', \uy' \ra) M_{\ell}(\ux')d \s
(\ux') = - \sigma_m \frac{k}{2(\l+k)} \delta_{k, \ell} M_{\ell}(\uy'),\\
\int_{\mS^{m-1}} (\ux' \wedge \uy' ) C_{k-1}^{\lambda+1}(\la\ux', \uy' \ra) \ux' M_{\ell}(\ux')d
\s (\ux') = \sigma_m \frac{k+2\l}{2(\l+k)} \delta_{k, \ell+1} \uy' M_{\ell}(\uy'),
\end{gather*}
where $C_{k}^{\lambda}(\la\ux',\uy'\ra)$ is the $k$-th Gegenbauer polynomial in the variable
$\la\ux',\uy'\ra$.
\end{lemma}

\begin{proof}
The f\/irst integral is trivial: $M_{\ell}$ is a~spherical harmonic of degree $\ell$ and
$C_{k}^{\lambda}(\la\ux',\uy'\ra)$ is the reproducing kernel for spherical harmonics of degree
$k$ (see, e.g.,~\cite{MR1827871}).
The second integral immediately follows, because $\ux'M_{\ell}(\ux')\in\cH_{\ell+1}$.

The other two integrals are a~bit more complicated.
We show how to obtain the last one.
First rewrite
$(\ux'\wedge\uy')\ux'=\uy'-\la\ux',\uy'\ra\ux'$.
The f\/irst term then follows from the f\/irst integral.
For the second term, we use the recursive property of Gegenbauer polynomials:
\begin{gather*}%\label{Geg2}
w C^{\l+1}_{n-1}(w)=\frac{n}{2(n+\l)}C^{\l+1}_{n}(w)+\frac{n+2\l}{2(n+\l)}
C^{\l+1}_{n-2}(w).
\end{gather*}
The result then follows by collecting everything.
\end{proof}

We can use this lemma to determine the reproducing kernels.
This is the subject of the following proposition.

\begin{proposition}
\label{reprkernels}
For $k\in\mN^{*}$ put
\begin{gather*}
P_{k}(\ux', \uy') = \frac{k+2\l}{2\l} C_{k}^{\l}(\la\ux', \uy' \ra) -(\ux' \wedge \uy' )
C_{k-1}^{\lambda+1}(\la\ux', \uy' \ra),\\
Q_{k-1}(\ux', \uy') = \frac{k}{2\l} C_{k}^{\l}(\la\ux', \uy' \ra) +(\ux' \wedge \uy' )
C_{k-1}^{\lambda+1}(\la\ux', \uy' \ra)
\end{gather*}
with $P_{0}(\ux',\uy')=C_{0}^{\l}(0)=1$.
Then
\begin{gather*}
\int_{\mS^{m-1}} P_{k}(\ux', \uy') M_{\ell}(\ux')d \s (\ux') = \sigma_m \delta_{k, \ell}
M_{\ell}(\uy'),\\
\int_{\mS^{m-1}} P_{k}(\ux', \uy') \ux'M_{\ell}(\ux')d \s (\ux') = 0
\end{gather*}
and
\begin{gather*}
\int_{\mS^{m-1}} Q_{k-1}(\ux', \uy') M_{\ell}(\ux')d \s (\ux') = 0,\\
\int_{\mS^{m-1}} Q_{k-1}(\ux', \uy') \ux'M_{\ell}(\ux')d \s (\ux') = \sigma_m \delta_{k, \ell+1}
\uy'M_{\ell}(\uy').
\end{gather*}
\end{proposition}

\begin{proof}
This follows immediately from Lemma~\ref{auxintegrals}.
\end{proof}

\begin{remark}
Note that, as expected, $P_{k}(\ux',\uy')+Q_{k-1}(\ux',\uy')=
\frac{\l+k}{\l}C_{k}^{\l}(\la\ux',\uy'\ra)$, which is the reproducing kernel for the space of
spherical harmonics of degree $k$.
\end{remark}

\begin{remark}
When the dimension $m=2$ and hence $\lambda=0$, the reproducing kernel is still well-def\/ined by
using the well-known relation~\cite[(4.7.8)]{Sz}
\begin{gather*}
\lim_{\lambda\rightarrow0}\lambda^{-1}C_k^\lambda(w)=(2/k)\cos k\theta,\qquad w=\cos
\theta,\qquad k\geq1.
\end{gather*}
\end{remark}

We will also need the following lemma.

\begin{lemma}
\label{ReprKernOrth}
The reproducing kernels satisfy the following properties, for all $k,l\in\mN$:
\begin{gather*}
\int_{\mS^{m-1}} P_{k}(\uy', \ux') P_{\ell}(\uz', \uy') d \s (\uy') = \sigma_m \delta_{k \ell
}P_{\ell}( \uz', \ux'),\\
\int_{\mS^{m-1}} P_{k}( \uy', \ux') Q_{\ell}( \uz', \uy') d \s (\uy') = 0,\\
\int_{\mS^{m-1}} Q_{k}( \uy', \ux') Q_{\ell}(\uz', \uy') d \s (\uy') = \sigma_m \delta_{k \ell }
Q_{\ell}(\uz', \ux').
\end{gather*}
\end{lemma}

\begin{proof}
This follows immediately using Lemma 7.6 and 7.10 from~\cite{DBXu}.
\end{proof}

\begin{remark}
Mind the order of the variables in the previous lemma.
The kernels $P_{k}(\ux',\uy')$ and $Q_{k}(\ux',\uy')$ are not symmetric.
\end{remark}

\section{The series representation of the holomorphic semigroup}\label{section6}

The aim of the present section is to investigate basic properties of the holomorphic semigroup
def\/ined by
\begin{gather*}
\cF_{\dD}^{\o}=e^{\o\left(\frac{1}{2}+\frac{\mu-1}{2(1+c)}\right)}e^{\frac{-\o}{
2(1+c)^{2}}\left(\dD^{2}-(1+c)^{2}\ux^{2}\right)},\qquad\operatorname{Re}\o\geq0,
\end{gather*}
acting on the space $\cL^{2}_{0,c}(\mR^{m})$.
We start with the following general statement.

\begin{theorem}\label{PropSemigroup}
Suppose $c>-1$.
Then
\begin{enumerate}\itemsep=0pt
\item[$1.$]
For any $t,\ell\in\mN$ and $m\in\{1,\ldots,\dim\cM_{\ell}\}$, the function
$\phi_{t,\ell,m}$ is an eigenfunction of the
operator $\cF_{\dD}^{\o}$:
\begin{gather*}
\cF_{\dD}^{\o}(\phi_{t,\ell,m})=e^{-\o t}e^{-\frac{\o\ell}{(1+c)}}\phi_{t,\ell,m}.
\end{gather*}
\item[$2.$] $\cF_{\dD}^{\o}$ is a~continuous operator on $\cL^{2}_{0,c}(\mR^{m})$ for all $\o$ with $\operatorname{Re}
\o\geq0$, in particular
\begin{gather*}
||\cF_{\dD}^{\o}(f)||\leq||f||
\end{gather*}
for all $f\in\cL^{2}_{0,c}(\mR^{m})$.
\item[$3.$]
If $\operatorname{Re}\o>0$, then $\cF_{\dD}^{\o}$ is a~Hilbert--Schmidt operator on $\cL^{2}_{0,c}(\mR^{m})$.
\item[$4.$]
If $\operatorname{Re}\o=0$, then $\cF_{\dD}^{\o}$ is a~unitary operator on $\cL^{2}_{0,c}(\mR^{m})$.
\end{enumerate}
\end{theorem}

\begin{proof}
$(1)$ is an immediate consequence of Theorem~\ref{HarmOsc}.
For $(2)$, let $f$ be an element in $\cL^{2}_{0,c}(\mR^{m})$ and expand it with respect to the
(normalized) basis $\{\phi_{t,\ell,m}\}$ as
\begin{gather*}
f=\sum_{t,\ell,m}a_{t,\ell,m}\phi_{t,\ell,m}.
\end{gather*}
Then one has, using orthogonality,
\begin{gather*}
||\cF_{\dD}^{\o}(f)||^2 = \sum_{t, \ell, m} |a_{t,\ell,m}|^{2} e^{-2 (\operatorname{Re}{\o}) t} e^{-\frac{2(\operatorname{Re}{\o}) \ell}{(1+c)}}
\leq \sum_{t, \ell, m} |a_{t,\ell,m}|^{2} = ||f||^2
\end{gather*}
because $\operatorname{Re}{\o}\geq0$.

As for $(3)$, we have to show that the Hilbert--Schmidt norm is f\/inite.
We compute
\begin{gather*}
||\cF_{\dD}^{\o}||_{\rm HS}^{2}
= \sum_{t, k, \ell} ||\cF_{\dD}^{\o}(\phi_{t,k,\ell})||^{2}
= \sum_{t, k,\ell} e^{-2 (\operatorname{Re}{\o}) t} e^{-\frac{2 (\operatorname{Re}{\o}) k}{(1+c)}}
= \sum_{t, k} e^{-2 (\operatorname{Re}{\o}) t} e^{-\frac{2 (\operatorname{Re}{\o}) k}{(1+c)}} \dim_{\mR}{\cM_{k}}\\
{} = \sum_{t, k} e^{-2 (\operatorname{Re}{\o}) t} e^{-\frac{2 (\operatorname{Re}{\o}) k}{(1+c)}} \frac{(k + m - 2)!}{k!(m - 2)!}2^{m}
= \sum_{t} e^{-2 (\operatorname{Re}{\o}) t} \sum_{k} e^{-\frac{2 (\operatorname{Re}{\o}) k}{(1+c)}} \frac{(k + m - 2)!}{k!(m - 2)!}2^{m}.
\end{gather*}
Using the ratio test, we see that these series are convergent for $\operatorname{Re}{\o}>0$.

$(4)$ follows immediately, because when $\operatorname{Re}{\o}=0$ the eigenvalues all have unit norm.
\end{proof}

We have already observed that $\cF_{\dD}^{\o}$ is a~Hilbert--Schmidt operator
for $\operatorname{Re}\o>0$ and a~unitary operator for $\operatorname{Re}\o=0$.
The Schwartz kernel theorem implies
that $\cF_{\dD}^{\o}$ can be expressed by a~distribution kernel
$K(x,y;\omega)$, so
\begin{gather*}
\big(\cF_{\dD}^{\o}f\big)(y)=\int_{\mR^m}K(x,y;\omega)f(x)h(r_{x})dx,
\end{gather*}
and $K(x,y;\omega)h(r_{x})$ is a~tempered distribution
on $\mR^m\times\mR^m$.

\subsection[The case $\operatorname{Re}\o>0$]{The case $\boldsymbol{\operatorname{Re}\o>0}$}

Using the reproducing kernels of Section~\ref{section5}, we can now make a~reasonable ansatz for
the kernel of the full holomorphic semigroup.
We want to write this semigroup as
\begin{gather*}
\cF_{0,c}^{\o}(f)(y)=\sigma_{m}^{-1}\int_{\mR^{m}}K(x,y;\omega)f(x)h(r_{x})dx
\end{gather*}
with $K(x,y;\omega)=K_{0}(x,y;\omega)+K_{1}(x,y;\omega)$ and
\begin{gather}
K_{0} = e^{-\frac{\coth{\o}}{2} (r^{2}+ s^{2})}\sum_{k=0}^{+\infty} \alpha_{k} z^{\frac{k}{1+c}}
\widetilde{J}_{\frac{\gamma_{k}}{2} -1}\left(\frac{iz}{\sinh{\o}}\right) P_{k}(\ux', \uy'),\nonumber\\
K_{1} = e^{-\frac{\coth{\o}}{2} (r^{2}+ s^{2})} \sum_{k=0}^{+\infty} \beta_{k} z^{1+\frac{k}{1+c}}
\widetilde{J}_{\frac{\gamma_{k}}{2}}\left(\frac{iz}{\sinh{\o}}\right) Q_{k}(\ux', \uy').\label{seriessemigroup}
\end{gather}
Here $r=|\ux|$, $s=|\uy|$ and $z=|\ux||\uy|$.
We also used the notation $\widetilde{J}_{\nu}(t)=(t/2)^{-\nu}J_{\nu}(t)$.
Now we determine the complex constants $\{\alpha_{k}\}$ and $\{\beta_{k}\}$ such that this
integral transform coincides with
\begin{gather*}
\cF_{\dD}^{\o}=e^{\o\left(\frac{1}{2}+\frac{\mu-1}{2(1+c)}\right)}e^{\frac{-\o}{
2(1+c)^{2}}\left(\dD^{2}-(1+c)^{2}\ux^{2}\right)}
\end{gather*}
on the basis $\{\phi_{t,\ell,m}\}$.

We calculate
\begin{gather*}
\sigma_{m}^{-1}\int_{\mR^{m}} K_{0}(x,y;\omega) \phi_{2t,\ell,m} (x)dx = \alpha_{\ell}
M_{\ell}^{(m)}(\uy') e^{-\frac{\coth{\o}}{2}s^{2}}s^{\frac{\ell}{1+c}}\left( \frac{i s}{2
\sinh{\o}}\right)^{-\gamma_{\ell}/2+1}\\
{}
\times \int_{0}^{+\infty} r^{\gamma_{\ell}/2} e^{-\frac{(\coth{\o}+1)}{2}r^{2}}
J_{\gamma_{\ell}/2 -1} \left( \frac{i r s}{\sinh{\o}}\right)L_{t}^{\gamma_{\ell}/2-1}(r^{2}) dr
=\alpha_{\ell} \frac{e^{-2\o t} 2^{\gamma_{\ell}/2-1}}{(\coth{\o}+1)^{\gamma_{\ell}/2}}
\phi_{2t,\ell,m}(y),
\end{gather*}
where we used the identity (see~\cite[Corollary 4.6]{Orsted2})
\begin{gather*}
2\int_{0}^{+\infty}r^{\alpha+1}J_{\a}(r\beta)L_{j}^{\a}(r^{2})e^{-\d r^{2}}dr=
\frac{(\d-1)^{j}\beta^{\alpha}}{2^{\alpha}\d^{\alpha+j+1}}L_{j}^{\a}\left(\frac{\beta^{2}}{4
\d(1-\d)}\right)e^{-\frac{\beta^{2}}{4\d}}.
\end{gather*}
Similarly, we f\/ind
\begin{gather*}
\sigma_{m}^{-1} \int_{\mR^{m}} K_{0}(x,y;\omega) \phi_{2t+1,\ell,m}(x) dx = 0,\qquad
\sigma_{m}^{-1} \int_{\mR^{m}} K_{1}(x,y;\omega) \phi_{2t,\ell,m} (x)dx = 0,\\
\sigma_{m}^{-1} \int_{\mR^{m}} K_{1}(x,y;\omega) \phi_{2t+1,\ell,m} (x) dx = \beta_{\ell}
\frac{e^{-2\o t} 2^{\gamma_{\ell}/2}}{(\coth{\o}+1)^{\gamma_{\ell}/2+1}} \phi_{2t+1,\ell,m}(y).
\end{gather*}
Hence we obtain by comparison with~\eqref{eigvalsFourier}
\begin{gather*}
\alpha_{\ell} = e^{-\frac{\o \ell}{(1+c)}}
\frac{(\coth{\o}+1)^{\gamma_{\ell}/2}}{2^{\gamma_{\ell}/2-1}} = 2 e^{\frac{\o \d}{2}} (2
\sinh{\o})^{-\gamma_{\ell}/2},\qquad
\beta_{\ell} = \frac{ \alpha_{\ell}}{2\sinh{\o}}.
\end{gather*}

We summarize our results in the following theorem.

\begin{theorem}
\label{seriesholom}
Let $\operatorname{Re}\o>0$ and $c>-1$.
Put
\begin{gather*}K(x,y;\omega)=e^{-\frac{\coth{\o}}{2}(r^{2}+s^{2})}\left(A(z,w)+\ux\wedge\uy
B(z,w)\right)
\end{gather*} with
\begin{gather*}
A(z,w) = \sum_{k=0}^{+\infty} \left(\alpha_{k} \frac{k+2\l}{2\l} z^{\frac{k}{1+c}}
\widetilde{J}_{\frac{\gamma_{k}}{2}-1}\left(\frac{iz}{\sinh{\o}}\right) +\frac{
\alpha_{k-1}}{4\sinh{\o}} \frac{k}{\l} z^{\frac{k+c}{1+c}}
\widetilde{J}_{\frac{\gamma_{k-1}}{2}}\left(\frac{iz}{\sinh{\o}}\right)\right) C_{k}^{\l}(w),\\
B(z,w) = \sum_{k=1}^{+\infty} \left(-\alpha_{k} z^{\frac{k}{1+c}-1}
\widetilde{J}_{\frac{\gamma_{k}}{2}-1}\left(\frac{iz}{\sinh{\o}}\right) +
\frac{\alpha_{k-1}}{2\sinh{\o}} z^{\frac{k+c}{1+c}-1}
\widetilde{J}_{\frac{\gamma_{k-1}}{2}}\left(\frac{iz}{\sinh{\o}}\right)\right) C_{k-1}^{\l+1}(w)
\end{gather*}
for $z=|\ux||\uy|$, $w=\la\ux,\uy\ra/z$, $\alpha_{-1}=0$ and
$%\begin{gather*}
\alpha_{k}=2e^{\frac{\o\d}{2}}(2\sinh{\o})^{-\gamma_{k}/2}.
$%\end{gather*}

Then these series are convergent and the integral transform defined on $\cL^{2}_{0,c}(\mR^{m})$ by
\begin{gather*}
\cF_{0,c}^{\o}(f)(y)=\sigma_{m}^{-1}\int_{\mR^{m}}K(x,y;\omega)f(x)h(r_{x})dx
\end{gather*}
coincides with the operator $\cF_{\dD}^{\o}=e^{\o\left(\frac{1}{2}+\frac{\mu-1}{2(1+c)}
\right)}e^{\frac{-\o}{2(1+c)^{2}}\left(\dD^{2}-(1+c)^{2}\ux^{2}\right)}$ on the basis $\{
\phi_{t,\ell,m}\}$.
\end{theorem}

\begin{proof}
We have already shown that the integral transform coincides with the operator $\cF_{\dD}^{\o}=
e^{\o\left(\frac{1}{2}+\frac{\mu-1}{2(1+c)}\right)}e^{\frac{-\o}{2(1+c)^{2}}\left(\dD^{2}-
(1+c)^{2}\ux^{2}\right)}$ on the basis $\{\phi_{t,\ell,m}\}$.
So we only have to show that the series are convergent.
We do this for the term
\begin{gather*}
\sum_{k=0}^{+\infty} (2 \sinh{\o})^{-\gamma_{k}/2} \frac{k+2\l}{2\l} z^{\frac{k}{1+c}}
\widetilde{J}_{\frac{\gamma_{k}}{2}-1}\left(\frac{iz}{\sinh{\o}}\right) C_{k}^{\l}(w)\\
\qquad
{}= (2 \sinh{\o})^{-\d/2} \sum_{k=0}^{+\infty} \frac{k+2\l}{2\l} \left(\frac{z}{2
\sinh{\o}}\right)^{\frac{k}{1+c}}
\widetilde{J}_{\frac{\gamma_{k}}{2}-1}\left(\frac{iz}{\sinh{\o}}\right) C_{k}^{\l}(w),
\end{gather*}
the other ones are treated in a~similar fashion.
We obtain
\begin{gather*}
\left| \sum_{k=0}^{+\infty} \frac{k+2\l}{2\l} \left(\frac{z}{2 \sinh{\o}}\right)^{\frac{k}{1+c}}
\widetilde{J}_{\frac{\gamma_{k}}{2}-1}\left(\frac{iz}{\sinh{\o}}\right) C_{k}^{\l}(w) \right| \\
\qquad
{}\leq \frac{B(\l)}{2}e^{\left|\operatorname{Im}{\frac{iz}{\sinh{\o}}}\right|}\sum_{k=0}^{+\infty}(k+2\l)\left|\frac{z}{2
\sinh{\o}}\right|^{\frac{k}{1+c}} \frac{1}{\Gamma(\gamma_{k}/2)} k^{2\l-1}
\end{gather*}
using formula~\eqref{estimateGeg} and \eqref{estimateBessel}.
As the term $\Gamma(\gamma_{k}/2)$ is dominant, the series clearly converges.
\end{proof}
\setcounter{section}{8}
\setcounter{subsection}{0}
\setcounter{subsubsection}{0}
\setcounter{equation}{1}
% Original source lines 1474--1515
\appendix

\section{Properties of Laguerre and Gegenbauer polynomials}\label{appendixA}

The generalized Laguerre polynomials $L_k^{(\alpha)}$ for $k\in\mN$ are def\/ined as
\begin{gather*}%\label{genLagdef}
L_k^{(\alpha)}(t)=\sum_{j=0}^{k}\frac{\Gamma(k+\alpha+1)}{j!(k-j)!\Gamma(j+\alpha+1)}(-t)^j
\end{gather*}
and satisfy the orthogonality relation (when $\alpha>-1$)
\begin{gather*}
\int_{0}^\infty t^{\alpha}
L^{(\alpha)}_k(t)L^{(\alpha)}_l(t)\exp(-t)dt=\delta_{kl}\frac{\Gamma(k+\alpha+1)}{k!}.
\end{gather*}

The Gegenbauer polynomials $C^{(\alpha)}_k(t)$ are a~special case of the Jacobi polynomials.
For $k\in\mN$ and $\alpha>-1/2$ they are def\/ined as
\begin{gather*}
C_k^{(\alpha)}(t)=\sum_{j=0}^{\lfloor
k/2\rfloor}(-1)^j\frac{\Gamma(k-j+\alpha)}{\Gamma(\alpha)j!(k-2j)!}(2t)^{k-2j}
%\label{GegenbauerCoeffs}
\end{gather*}
and satisfy the orthogonality relation
\begin{gather*}
\int_{-1}^1C_k^{(\alpha)}(t)C_l^{(\alpha)}(t)\big(1-t^2\big)^{\alpha-\frac{1}{2}}dt=\delta_{kl}\frac{\pi2^{1
-2\alpha}\Gamma(k+2\alpha)}{k!(k+\alpha)(\Gamma(\alpha))^2}.
\end{gather*}
One can prove that there exists a~constant $B(\a)$ such that
\begin{gather}\label{estimateGeg}
\sup_{-1\leq t\leq1}\left|\frac{1}{\a}C_k^{(\alpha)}(t)\right|\leq B(\a)k^{2\a-1},\qquad
\forall\, k\in\mN,
\end{gather}
see~\cite[Lemma 4.9]{Orsted2}.

The Bessel function $J_{\nu}(z)$ is def\/ined using the following Taylor series
\begin{gather*}
J_{\nu}(z)=\sum_{k=0}^\infty\frac{(-1)^k}{k!\Gamma(k+\nu+1)}
{\left({\frac{z}{2}}\right)}^{2k+\nu}.
\end{gather*}
For $z\in\mC$ and $\nu\geq-1/2$ one has the inequality (see, e.g.,~\cite{Sz})
\begin{gather}\label{estimateBessel}
\left|\left(\frac{z}{2}\right)^{-\nu}J_{\nu}(z)\right|\leq\frac{1}{\Gamma(\nu+1)}e^{|\operatorname{Im}{z}|}.
\end{gather}
\begin{thebibliography}{99}
\bibitem[1]{BCT}
Barbasch D., Ciubotaru D., Trapa P.E., Dirac cohomology for graded af\/f\/ine
{H}ecke algebras, \href{http://dx.doi.org/10.1007/s11511-012-0085-3}{\textit{Acta Math.}}
\textbf{209} (2012), 197--227,
\href{http://arxiv.org/abs/1006.3822}{arXiv:1006.3822}.

\bibitem[2]{Orsted2}
Ben~Sa{\"{\i}}d S., Kobayashi T., {\O}rsted B., Laguerre semigroup and {D}unkl
operators, \href{http://dx.doi.org/10.1112/S0010437X11007445}{\textit{Compos.
Math.}} \textbf{148} (2012), 1265--1336,
\href{http://arxiv.org/abs/0907.3749}{arXiv:0907.3749}.

\bibitem[3]{MR2269930}
Ca{\c{c}}{\~a}o I., Constales D., Krausshar R.S., On the role of arbitrary
order {B}essel functions in higher dimensional {D}irac type equations,
\href{http://dx.doi.org/10.1007/s00013-006-1791-x}{\textit{Arch.
Math.
(Basel)}} \textbf{87} (2006), 468--477.

\bibitem[5]{DBGeg}
De~Bie H., De~Schepper N., Clif\/ford--{G}egenbauer polynomials related to the
{D}unkl {D}irac operator, \textit{Bull.
Belg.
Math.
Soc.
Simon Stevin}
\textbf{18} (2011), 193--214.

\bibitem[7]{H12}
De~Bie H., {\O}rsted B., Somberg P., Sou{\v{c}}ek V., Dunkl operators and a
family of realizations of {$\mathfrak{osp}(1\vert2)$},
\href{http://dx.doi.org/10.1090/S0002-9947-2012-05608-X}{\textit{Trans.
Amer.
Math.
Soc.}} \textbf{364} (2012), 3875--3902, \href{http://arxiv.org/abs/0911.4725}{arXiv:0911.4725}.

\bibitem[8]{DBXu}
De~Bie H., Xu Y., On the {C}lif\/ford--{F}ourier transform,
\href{http://dx.doi.org/10.1093/imrn/rnq288}{\textit{Int.
Math.
Res.
Not.}} \textbf{2011} (2011), no.~22, 5123--5163,
\href{http://arxiv.org/abs/1003.0689}{arXiv:1003.0689}.

\bibitem[9]{MR1169463}
Delanghe R., Sommen F., Sou{\v{c}}ek V., Clif\/ford algebra and spinor-valued
functions.
A~function theory for the Dirac operator,
\href{http://dx.doi.org/10.1007/978-94-011-2922-0}{\textit{Mathematics and
its Applications}}, Vol.~53, Kluwer Academic Publishers Group, Dordrecht,
1992.

\bibitem[10]{MR951883}
Dunkl C.F., Dif\/ferential-dif\/ference operators associated to ref\/lection groups,
\href{http://dx.doi.org/10.2307/2001022}{\textit{Trans.
Amer.
Math.
Soc.}} \textbf{311} (1989), 167--183.

\bibitem[11]{Du4}
Dunkl C.F., Hankel transforms associated to f\/inite ref\/lection groups, in
Hypergeometric Functions on Domains of Positivity, {J}ack Polynomials, and
Applications ({T}ampa, {FL}, 1991),
\href{http://dx.doi.org/10.1090/conm/138/1199124}{\textit{Contemp.
Math.}}, Vol.~138, Amer.
Math.
Soc., Providence, RI, 1992, 123--138.

\bibitem[12]{MR1273532}
Dunkl C.F., de~Jeu M.F.E., Opdam E.M., Singular polynomials for f\/inite
ref\/lection groups, \href{http://dx.doi.org/10.2307/2154950}{\textit{Trans.
Amer.
Math.
Soc.}} \textbf{346} (1994),
237--256.

\bibitem[13]{MR1827871}
Dunkl C.F., Xu Y., Orthogonal polynomials of several variables,
\href{http://dx.doi.org/10.1017/CBO9780511565717}{\textit{Encyclopedia of Mathematics and its
Applications}}, Vol.~81, Cambridge
University Press, Cambridge, 2001.

\bibitem[15]{MR1773773}
Frappat L., Sciarrino A., Sorba P., Dictionary on {L}ie algebras and
superalgebras, Academic Press Inc., San Diego, CA, 2000.

\bibitem[16]{MR1130821}
Gilbert J.E., Murray M.A.M., Clif\/ford algebras and {D}irac operators in
harmonic analysis, \href{http://dx.doi.org/10.1017/CBO9780511611582}{\textit{Cambridge Studies in
Advanced Mathematics}},
Vol.~26, Cambridge University Press, Cambridge, 1991.

\bibitem[17]{MR0974332}
Howe R., The oscillator semigroup, in The Mathematical Heritage of {H}ermann
{W}eyl ({D}urham, {NC}, 1987), \textit{Proc.
Sympos.
Pure Math.}, Vol.~48,
Amer.
Math.
Soc., Providence, RI, 1988, 61--132.

\bibitem[18]{Humph}
Humphreys J.E., Ref\/lection groups and {C}oxeter groups, \textit{Cambridge
Studies in Advanced Mathematics}, Vol.~29, Cambridge University Press,
Cambridge, 1990.

\bibitem[23]{Orsted}
{\O}rsted B., Somberg P., Sou{\v{c}}ek V., The {H}owe duality for the {D}unkl
version of the {D}irac operator, \href{http://dx.doi.org/10.1007/s00006-009-0166-3}{\textit{Adv.
Appl.
Clif\/ford Algebr.}}
\textbf{19} (2009), 403--415.

\bibitem[25]{MR2022853}
R{\"o}sler M., Dunkl operators: theory and applications, in Orthogonal
Polynomials and Special Functions ({L}euven, 2002),
\href{http://dx.doi.org/10.1007/3-540-44945-0_3}{\textit{Lecture Notes in
Math.}}, Vol.~1817, Springer, Berlin, 2003, 93--135,
\href{http://arxiv.org/abs/math.CA/0210366}{math.CA/0210366}.

\bibitem[27]{Sz}
Szeg{\"{o}} G., Orthogonal polynomials, \textit{American Mathematical Society,
Colloquium Publications}, Vol.~23, 4th~ed., American Mathematical Society,
Providence, R.I., 1975.

\end{thebibliography}
\end{document}
